% =====================================================================
% §1.3 — Asignación de especies por contenedor
% =====================================================================
\section{Asignación de especies por contenedor}
\label{sec:especies}

\textbf{Nota de revisión.} La selección final de especies y sus
\emph{setpoints} óptimos están sujetos a validación tras el curso de
cultivo del operador. Esta sección consolida valores de partida
documentados en literatura para tres especies de uso comercial común.
La tabla \ref{tab:especies} sirve como base para los \emph{setpoints}
del firmware (Fase~2).

% ---------------------------------------------------------------------
\subsection{Parámetros ambientales por especie (literatura)}
\label{sec:especies:tabla}

\begin{table}[H]
  \centering
  \caption{Rangos ambientales óptimos de fructificación reportados en
    literatura. Adaptado de Stamets~\cite{stamets2000} y Chang \&
    Miles~\cite{chang2004}.}
  \label{tab:especies}
  \begin{tabularx}{\linewidth}{Xccccc}
    \toprule
    \textbf{Especie} & T [\si{\celsius}] & HR [\%] & CO$_2$ [ppm] & Lux & Fotop. [h] \\
    \midrule
    \paramEspecieA & 18--24 & 85--95 & $<$1000 & 500--1500 & 12 \\
    \paramEspecieB & 18--24 & 90--95 & $<$800  & 500--1000 & 12 \\
    \paramEspecieC & 12--18 & 80--90 & $<$1000 & 500--2000 & 12 \\
    \bottomrule
  \end{tabularx}
\end{table}

\paragraph{Notas operativas.}
\begin{itemize}[leftmargin=*,itemsep=0pt]
  \item \textbf{Oyster} (\paramEspecieA): especie más permisiva y rápida.
        Ideal como especie de entrenamiento operativo. Tolera CO$_2$
        más alto a costa de elongación del estípite.
  \item \textbf{Lion's mane} (\paramEspecieB): sensible a corrientes
        directas de aire. Forma "barbas" con HR sostenida $>$\,90\,\%;
        bajo deshidratación produce ramificación corta.
  \item \textbf{Shiitake} (\paramEspecieC): prefiere temperaturas más
        bajas y choque térmico inicial (cold shock $\approx 10$\,\si{\celsius}
        durante 24--48\,h) para inducir pinning.
\end{itemize}

\subsection{Asignación propuesta (tentativa)}
\label{sec:especies:asignacion}

Justificación de la asignación contenedor$\leftrightarrow$especie
basada en la jerarquía térmica natural del rack (calor asciende):

\begin{table}[H]
  \centering
  \caption{Asignación propuesta de especies por nivel.}
  \label{tab:asignacion}
  \begin{tabularx}{\linewidth}{cXl}
    \toprule
    \textbf{Nivel} & \textbf{Especie} & \textbf{Justificación} \\
    \midrule
    1 (sup.) & \paramEspecieB & T más alta naturalmente; banda 18--24\,\si{\celsius}. \\
    2        & \paramEspecieA & Banda 18--24\,\si{\celsius}; alta rotación, acceso medio. \\
    3        & \paramEspecieC & Banda térmica más baja del rack (12--18\,\si{\celsius}). \\
    4 (inf.) & --- & Servicios (electrónica, fuentes, broker MQTT). \\
    \bottomrule
  \end{tabularx}
\end{table}

\pendientecurso{Validar la jerarquía térmica del rack midiendo en sitio
con termómetro IR antes de fijar la asignación. Reasignar si la
estratificación real difiere de la asumida.}

\subsection{Riesgo de contaminación cruzada}
\label{sec:especies:contaminacion}

La segregación física por contenedor mitiga contaminación cruzada
entre especies; sin embargo, durante el manejo (apertura de tapa,
cosecha) se generan corrientes que pueden transportar esporas. Medidas
operativas:

\begin{itemize}[leftmargin=*,itemsep=0pt]
  \item Abrir un contenedor a la vez.
  \item Manejar primero la especie con mayor sensibilidad
        (lion's mane) y al final la más robusta (oyster).
  \item Lavarse manos y desinfectar superficie de trabajo entre
        manipulaciones.
\end{itemize}

\subsection{Criterios de aceptación}
\label{sec:especies:aceptacion}

\begin{itemize}[leftmargin=*,itemsep=0pt]
  \item Tabla de parámetros respaldada por al menos dos referencias
        bibliográficas verificables.
  \item Asignación contenedor$\leftrightarrow$especie revalidada con
        mapeo térmico en sitio antes de cargar el primer lote.
  \item Setpoints de Fase~2 derivan directamente de
        tabla~\ref{tab:especies}.
\end{itemize}
